%=====================================================================
\chapter{Exploratory Data Analysis and Visualisation}\label{ch:eda}
%=====================================================================
\locator{2}{Part II --- From Data to Insight}

\begin{outcomes}
By the end of this chapter you will be able to:
\begin{tight}
  \item \textbf{Conduct} a structured EDA: univariate, then bivariate, then multivariate.
  \item \textbf{Select} the correct chart for a given pair of variable types.
  \item \textbf{Read} a correlation heatmap and recognise multicollinearity.
  \item \textbf{Identify} the common visualisation anti-patterns and explain the distortion each causes.
  \item \textbf{Design} a dashboard around the decision it supports rather than the data available.
\end{tight}
\end{outcomes}

\begin{prereq}
\chref{descriptive} (centre, spread, shape) and \chref{wrangling} (you cannot
explore data you have not cleaned).
\end{prereq}

\begin{keyterms}
exploratory data analysis $\bullet$ univariate, bivariate, multivariate $\bullet$
histogram $\bullet$ boxplot $\bullet$ scatterplot $\bullet$ correlation matrix
$\bullet$ heatmap $\bullet$ small multiples $\bullet$ multicollinearity $\bullet$
Simpson's paradox $\bullet$ data-ink $\bullet$ dashboard
\end{keyterms}

\section{What EDA Is For}

\begin{definitionbox}[Exploratory Data Analysis]
The systematic use of summaries and graphics to understand a dataset's structure,
find its problems and generate hypotheses --- \emph{before} any model is fitted.
Its output is not a chart; it is a list of things you now know and things you now
suspect.
\end{definitionbox}

EDA answers four questions, in order:
\begin{tightnum}
  \item What does each variable look like on its own? (distribution, outliers, missingness)
  \item How do pairs of variables relate? (correlation, group differences)
  \item Does anything change when a third variable is held constant? (confounding)
  \item Is anything here that should not be? (impossible values, leaks, duplicates)
\end{tightnum}

\section{Choosing the Right Chart}

The choice is almost entirely determined by the \emph{types} of the variables
involved --- which is why \chref{data} came first.

\dsfig{%
\begin{tikzpicture}[font=\scriptsize,
  cell/.style={rounded corners=2pt, draw=#1, fill=#1!8, text=#1!50!black,
               text width=3.0cm, align=center, minimum height=1.25cm, font=\scriptsize},
  hd/.style={rounded corners=2pt, fill=primary, text=white,
             text width=3.0cm, align=center, minimum height=0.7cm, font=\scriptsize\bfseries},
  rh/.style={rounded corners=2pt, fill=primary!75, text=white,
             text width=2.75cm, align=center, minimum height=1.25cm, font=\scriptsize\bfseries}]

\node[hd] at (2.6,0)  {ONE variable};
\node[hd] at (5.9,0)  {TWO: num vs num};
\node[hd] at (9.2,0)  {TWO: num vs category};
\node[hd] at (12.5,0) {MANY variables};

\node[rh] at (-0.55,-1.25) {NUMERIC};
\node[cell=secondary] at (2.6,-1.25)  {\textbf{Histogram}\\shape, skew, modes\\\emph{+ boxplot for outliers}};
\node[cell=greenhl]   at (5.9,-1.25)  {\textbf{Scatterplot}\\form, strength,\\direction, outliers};
\node[cell=goldhl]    at (9.2,-1.25)  {\textbf{Boxplot by group}\\or violin plot\\\emph{compare distributions}};
\node[cell=purplehl]  at (12.5,-1.25) {\textbf{Correlation heatmap}\\or pair plot};

\node[rh] at (-0.55,-2.85) {CATEGORICAL};
\node[cell=secondary] at (2.6,-2.85)  {\textbf{Bar chart}\\of counts\\\emph{never a pie chart}};
\node[cell=greenhl]   at (5.9,-2.85)  {\textbf{Grouped bar}\\or mosaic plot};
\node[cell=goldhl]    at (9.2,-2.85)  {\textbf{Stacked / grouped bar}\\proportions per group};
\node[cell=purplehl]  at (12.5,-2.85) {\textbf{Small multiples}\\one panel per level};

\node[rh] at (-0.55,-4.45) {OVER TIME};
\node[cell=secondary] at (2.6,-4.45)  {\textbf{Line chart}\\trend and seasonality};
\node[cell=greenhl]   at (5.9,-4.45)  {\textbf{Dual line}\\(shared axis only if\\units match)};
\node[cell=goldhl]    at (9.2,-4.45)  {\textbf{Line per group}\\max $\sim$5 lines};
\node[cell=purplehl]  at (12.5,-4.45) {\textbf{Heatmap}\\time $\times$ entity};
\end{tikzpicture}}%
{Chart selection by variable type. Almost every ``which chart should I use?'' question is
answered by identifying the types involved and reading off this grid.}{chart-selection}

\begin{alertbox}[Why Not a Pie Chart]
Humans compare lengths accurately and angles badly. A bar chart lets a reader rank
six categories at a glance; a pie chart of the same data does not. Reserve pies
for the one case they handle well: two or three parts of an obvious whole.
\end{alertbox}

\section{Bivariate Exploration}

\dsfig{%
\begin{tikzpicture}
\begin{groupplot}[
  group style={group size=4 by 1, horizontal sep=0.62cm},
  width=4.35cm, height=3.6cm,
  axis lines=left, tick label style={font=\tiny},
  title style={font=\tiny\bfseries}, xtick=\empty, ytick=\empty,
  xmin=0, xmax=10, ymin=0, ymax=10,
]
\nextgroupplot[title={\color{greenhl}strong positive, linear}]
\addplot[only marks, mark=*, mark size=1pt, greenhl!80] coordinates {
 (1,1.6)(1.7,2.1)(2.3,3.4)(3,3.2)(3.6,4.6)(4.2,4.2)(5,5.5)(5.6,5.2)
 (6.3,6.8)(7,6.4)(7.6,7.9)(8.3,7.6)(9,8.8)(2.6,2.6)(6.8,7.4)};
\addplot[primary, line width=0.8pt, domain=0.8:9.2] {0.92*x+0.6};

\nextgroupplot[title={\color{accent}strong, but not linear}]
\addplot[only marks, mark=*, mark size=1pt, accent!80] coordinates {
 (1,8.6)(1.7,6.4)(2.3,4.8)(3,3.6)(3.6,2.6)(4.2,2.0)(5,1.6)(5.6,1.8)
 (6.3,2.4)(7,3.4)(7.6,4.8)(8.3,6.4)(9,8.4)(2.6,4.0)(6.8,3.0)};
\node[font=\tiny, accent!80!black, align=center] at (axis cs:5,9.0) {$r \approx 0$ !};

\nextgroupplot[title={\color{secondary}no relationship}]
\addplot[only marks, mark=*, mark size=1pt, secondary!80] coordinates {
 (1,5.2)(1.7,2.4)(2.3,7.8)(3,4.1)(3.6,8.6)(4.2,2.2)(5,6.4)(5.6,3.2)
 (6.3,7.2)(7,4.6)(7.6,1.8)(8.3,6.0)(9,3.6)(2.6,5.8)(6.8,8.2)};

\nextgroupplot[title={\color{purplehl}one point drives $r$}]
\addplot[only marks, mark=*, mark size=1pt, purplehl!80] coordinates {
 (1,4.6)(1.5,5.2)(2,4.2)(2.4,5.6)(2.8,4.8)(3.2,5.4)(3.6,4.4)(4,5.0)(4.4,4.6)};
\addplot[only marks, mark=*, mark size=2.2pt, accent] coordinates {(9,9.2)};
\addplot[primary, line width=0.8pt, domain=0.8:9.2] {0.62*x+3.3};
\node[font=\tiny, accent, anchor=east] at (axis cs:8.6,9.2) {leverage};
\end{groupplot}
\end{tikzpicture}}%
{Four scatterplots. The second has a perfect relationship and a correlation near zero; the
fourth has a correlation near 0.9 created entirely by one point. This is why you plot the
data before trusting any correlation coefficient.}{scatter-patterns}

\begin{definitionbox}[Pearson Correlation]
$r$ measures the strength and direction of a \emph{linear} relationship, on a scale
from $-1$ to $+1$. Zero means no \emph{linear} relationship --- it does not mean no
relationship, as \dsref{scatter-patterns} shows.
\end{definitionbox}

\begin{worked}[Reading a Correlation Matrix]
An engagement dataset gives these correlations with the target
\texttt{final\_mark}:

\smallskip
\begin{center}
\begin{tabular}{@{}lrr@{}}
\toprule
\textbf{Feature} & \textbf{$r$ with mark} & \textbf{$r$ with attendance} \\
\midrule
attendance\_pct     & 0.62 & 1.00 \\
lectures\_watched   & 0.58 & 0.94 \\
quizzes\_done       & 0.55 & 0.41 \\
forum\_posts        & 0.21 & 0.18 \\
days\_since\_login  & $-0.49$ & $-0.73$ \\
\bottomrule
\end{tabular}
\end{center}

\smallskip
\textbf{Reading 1 --- predictive strength.} Attendance is the strongest single
predictor at $r = 0.62$. Squaring it, $r^2 = 0.38$: attendance alone accounts for
about 38\% of the variation in marks. Useful, but 62\% is elsewhere.

\textbf{Reading 2 --- multicollinearity.} \texttt{attendance\_pct} and
\texttt{lectures\_watched} correlate at 0.94 with \emph{each other}. They are
nearly the same variable measured twice. Including both adds almost no information
and makes the coefficients of a linear model unstable and uninterpretable
(\chref{regression}). Keep one.

\textbf{Reading 3 --- the weak feature.} \texttt{forum\_posts} at $r = 0.21$ looks
negligible. Do not drop it yet: correlation only sees linear effects, and a
feature that is weak alone can be valuable in combination (\chref{features}).

\textbf{Reading 4 --- sign.} \texttt{days\_since\_login} is negative, as it should
be. A positive sign here would have signalled a data error worth chasing.
\end{worked}

\section{The Third Variable}

\begin{alertbox}[Simpson's Paradox]
A relationship visible in aggregated data can \emph{reverse} when the data is
split by a third variable. A university may show that students who use the online
tutorials score \emph{lower} overall --- while within every individual programme,
tutorial users score higher. The aggregate is dominated by the fact that harder
programmes both recommend tutorials more and mark more severely.

The practical instruction: whenever you show a relationship, ask what the obvious
grouping variable is, and re-plot within groups before you present it.
\chref{causality} treats this properly as confounding.
\end{alertbox}

\dsfig{%
\begin{tikzpicture}
\begin{groupplot}[
  group style={group size=2 by 1, horizontal sep=1.4cm},
  width=6.6cm, height=4.6cm,
  axis lines=left, tick label style={font=\tiny}, label style={font=\tiny},
  title style={font=\scriptsize\bfseries},
  xmin=0, xmax=10, ymin=25, ymax=85,
  xlabel={tutorial hours used}, ylabel={final mark},
]
\nextgroupplot[title={\color{accent}Aggregated: looks harmful}]
\addplot[only marks, mark=*, mark size=1.1pt, secondary!70] coordinates {
 (0.6,72)(1.1,74)(1.5,70)(1.9,76)(2.4,73)(2.8,78)(3.2,75)(1.3,71)(2.1,77)(0.9,69)};
\addplot[only marks, mark=triangle*, mark size=1.3pt, goldhl] coordinates {
 (6.2,48)(6.8,52)(7.3,46)(7.9,54)(8.4,50)(8.9,56)(9.3,52)(7.1,49)(8.1,53)(6.5,45)};
\addplot[accent, line width=1.2pt, domain=0.4:9.5] {77 - 3.4*x};
\node[font=\tiny, accent, anchor=north east] at (axis cs:9.6,80) {overall trend: negative};

\nextgroupplot[title={\color{greenhl}Split by programme: helpful in both}]
\addplot[only marks, mark=*, mark size=1.1pt, secondary!70] coordinates {
 (0.6,72)(1.1,74)(1.5,70)(1.9,76)(2.4,73)(2.8,78)(3.2,75)(1.3,71)(2.1,77)(0.9,69)};
\addplot[secondary, line width=1.2pt, domain=0.4:3.4] {68.5 + 2.4*x};
\addplot[only marks, mark=triangle*, mark size=1.3pt, goldhl] coordinates {
 (6.2,48)(6.8,52)(7.3,46)(7.9,54)(8.4,50)(8.9,56)(9.3,52)(7.1,49)(8.1,53)(6.5,45)};
\addplot[goldhl, line width=1.2pt, domain=6.0:9.5] {32 + 2.3*x};
\node[font=\tiny, secondary!60!black, anchor=west] at (axis cs:3.4,80) {Programme A};
\node[font=\tiny, goldhl!50!black, anchor=west] at (axis cs:5.0,36) {Programme B};
\end{groupplot}
\end{tikzpicture}}%
{Simpson's paradox. The same points, plotted twice. Aggregated, tutorial use appears to
lower marks; within each programme it raises them. Programme is a confounder --- students
on the harder programme use tutorials more \emph{and} score lower for unrelated
reasons.}{simpsons}

\section{Visualisation Anti-Patterns}

\rowcolors{2}{white}{lightbg}
\begin{longtable}{@{}L{4.2cm}L{5.2cm}L{5.1cm}@{}}
\toprule
\rowcolor{primary!15}
\textbf{Anti-pattern} & \textbf{The distortion} & \textbf{The fix} \\
\midrule
\endhead
Truncated $y$-axis on a bar chart & Bar length no longer encodes value; a 2\% change looks like a doubling & Bars must start at zero. Use a line chart if you need to zoom \\
Dual $y$-axes & The apparent correlation depends entirely on where you set the two scales & Two stacked panels sharing an $x$-axis \\
3-D effects & Perspective makes near bars look larger & Never use them \\
Rainbow colour scale & Not perceptually uniform; invents boundaries that are not in the data & Sequential (light-to-dark) for magnitude, diverging for signed \\
Red/green as the only cue & Invisible to about 1 in 12 men & Add shape or position; use blue/orange \\
Too many lines & Above five series nothing is legible & Small multiples, or highlight one and grey the rest \\
Averaging away the shape & A single mean hides bimodality entirely & Show the distribution: histogram, boxplot or strip plot \\
\bottomrule
\end{longtable}

\begin{conceptbox}[One Test for Any Chart]
Cover the title and hand it to a colleague. If they cannot state the main message
in one sentence within about ten seconds, the chart is not finished. This test
catches more problems than any style guide.
\end{conceptbox}

%---------------------------------------------------------------------
\begin{fourlenses}{Exploratory Analysis}
\lensLA{Plot engagement over the weeks of the semester as one line per student,
faintly, with the cohort median in bold --- the ``spaghetti plot''. The
\emph{shape} of individual disengagement (a cliff in week 5 versus a slow decline)
is invisible in any average and is exactly what the intervention design needs.
Always split by programme before presenting: Simpson's paradox is endemic in
education data.}

\lensPM{Plot estimated against actual effort as a scatterplot with the $y=x$ line
drawn. The systematic gap above that line \emph{is} the estimation bias, and it is
far more persuasive to a team than any summary statistic. A cumulative flow diagram
shows where work piles up; a burn-down chart alone does not.}

\lensIOT{With 400 devices you cannot plot 400 lines. Use a heatmap of device
$\times$ time with colour as the reading: a failing sensor appears as a stripe, a
site-wide outage as a vertical band. This one chart replaces 400 line charts and
reveals patterns none of them would show individually.}

\lensSEC{Log-scale nearly everything --- counts of events per source are long-tailed,
and a linear axis renders all but the largest invisible. Plot volume by hour of day
to establish the normal diurnal rhythm, because \emph{deviation from the rhythm} is
more informative than raw volume. Never truncate an axis on a security chart: it is
how a small anomaly is made to look like a crisis, and how a real one gets missed.}
\end{fourlenses}

%---------------------------------------------------------------------
\begin{lab}[A Standard EDA Pass]
\begin{lstlisting}[language=Python]
import pandas as pd, matplotlib.pyplot as plt, seaborn as sns

df = pd.read_csv("students.csv")

# --- 1. univariate: shape of every numeric column ---------------------
df.select_dtypes("number").hist(bins=30, figsize=(11, 7))
plt.tight_layout(); plt.show()

# --- 2. bivariate: relationships and multicollinearity ----------------
corr = df.select_dtypes("number").corr()
sns.heatmap(corr, annot=True, fmt=".2f", cmap="RdBu_r",
            center=0, vmin=-1, vmax=1)   # diverging scale for signed values
plt.show()

# Flag pairs that are nearly the same variable measured twice.
import numpy as np
high = (corr.abs() > 0.85) & (corr.abs() < 1.0)
print("multicollinear pairs:",
      [(a, b) for a in corr.index for b in corr.columns
       if high.loc[a, b] and a < b])

# --- 3. the third variable: always check before presenting ------------
sns.lmplot(data=df, x="tutorial_hours", y="final_mark",
           hue="programme", height=4)    # per-group, not aggregated
plt.show()
\end{lstlisting}

\textbf{R equivalent}, matching the repo's existing \texttt{mtcars} EDA lab:
\begin{lstlisting}[language=R]
library(tidyverse); library(GGally)
ggpairs(select(df, where(is.numeric)))
ggplot(df, aes(tutorial_hours, final_mark, colour = programme)) +
  geom_point() + geom_smooth(method = "lm")
\end{lstlisting}
\end{lab}

%---------------------------------------------------------------------
\begin{checkpoint}
\begin{tightnum}
  \item Two features correlate at $r = 0.96$ with each other. What is this called
        and what should you do before fitting a linear model?
  \item A scatterplot shows a clear U-shape and $r = 0.02$. Is there a relationship?
  \item You must compare failure rates across 40 servers over 30 days. Which chart,
        and why not 40 line charts?
\end{tightnum}
\end{checkpoint}

%---------------------------------------------------------------------
\begin{chaptersummary}
\begin{tight}
  \item EDA runs univariate $\rightarrow$ bivariate $\rightarrow$ multivariate, and
        its output is knowledge and hypotheses, not pictures.
  \item Chart choice follows from variable types; the grid in
        \dsref{chart-selection} answers almost every case.
  \item Pearson's $r$ measures \emph{linear} association only, and is easily
        dominated by a single high-leverage point. Always plot before trusting it.
  \item Correlations above about 0.85 between features signal multicollinearity;
        keep one of the pair.
  \item Simpson's paradox means an aggregate relationship can reverse within
        groups. Always check the obvious grouping variable.
  \item Truncated axes, dual axes, 3-D effects and rainbow scales are distortions,
        not styles.
\end{tight}
\end{chaptersummary}

\begin{reviewq}
  \item \textbf{(MCQ)} Pearson's $r$ of 0 means: (a)~no relationship (b)~no linear
        relationship (c)~the variables are independent (d)~the data is normal.
  \item \textbf{(MCQ)} For comparing the distribution of a numeric variable across
        five categories, the best chart is: (a)~pie chart (b)~stacked bar
        (c)~boxplot by group (d)~3-D column chart.
  \item \textbf{(Short)} Explain why a bar chart must start at zero but a line
        chart need not.
  \item \textbf{(Short)} Define multicollinearity and state one concrete problem it
        causes.
  \item \textbf{(Short)} Describe Simpson's paradox in two sentences and give an
        example from any domain.
  \item \textbf{(Applied)} You are handed a dataset of 15{,}000 login events with
        columns \texttt{user}, \texttt{timestamp}, \texttt{source\_ip},
        \texttt{success}, \texttt{duration\_s}. Write the EDA plan: list the charts
        you would produce, in order, and state what each is meant to reveal.
  \item \textbf{(Applied)} A manager presents a chart with a truncated $y$-axis
        showing incidents ``up 300\%'' (from 4 to 12). Rewrite the chart
        specification honestly, and say what additional context the reader needs.
\end{reviewq}
